% Lecture 6 notes: \input by notes.tex, \subimport'ed by the semester book.
% Exam questions (зачёт): 2 (бинарный алгоритм и его модификации),
% 4 (возведение в степень по модулю методом Монтгомери).
\begin{lecture}{6}{Быстрое возведение в степень}

\section{Зачем быстро возводить в степень}

Почти все «законные» операции асимметричной криптографии сводятся к возведению в степень.
\begin{itemize}
  \item \emph{Протокол Диффи"--~Хеллмана} (лекция~1 и лекция~1 второго семестра): Алиса
    и Боб вычисляют $g^a \bmod p$, $g^b \bmod p$ и общий ключ $(g^b)^a = (g^a)^b \bmod p$,
    где $p$ "--- простое число длиной $2048$--$3072$ бит, а показатели $a, b$ имеют длину
    не меньше $256$ бит.
  \item \emph{RSA} (лекция~5 второго семестра): шифрование $c = m^e \bmod N$, расшифрование
    $m = c^d \bmod N$; модуль $N$ и секретный показатель $d$ имеют длину $2048$ бит и больше.
  \item \emph{Тесты простоты} (лекция~8): тест Ферма и тест Миллера"--~Рабина проверяют
    сравнения вида $a^{n-1} \equiv 1 \pmod n$.
  \item \emph{Криптография на эллиптических кривых} (лекции~8--9 второго семестра): вычисление
    кратной точки $kP$ "--- то же возведение в степень, только в аддитивной записи.
\end{itemize}
Наивное вычисление $a^e = a \cdot a \cdots a$ требует $e - 1$ умножений, а при $e \approx 2^{2048}$
это число несравнимо больше числа атомов во Вселенной. В этой лекции мы увидим, что достаточно
$\bigO(\log e)$ умножений (точнее, около $1{,}5\log_2 e$, а с предвычислениями "--- $\log_2 e\,(1 + o(1))$),
и разберём, как выполнять эти умножения по модулю без делений (метод Монтгомери).

\section{Степени в моноиде}

Алгоритмы этой лекции работают в любом \emph{моноиде} "--- множестве $M$ с ассоциативной
операцией (будем записывать её как умножение) и нейтральным элементом $1$. Основные примеры:
кольцо вычетов $\Zn{m}$ по умножению, группа $\Znx{m}$, кольцо матриц $n \times n$, кольцо
многочленов по модулю фиксированного многочлена, группа точек эллиптической кривой.

\begin{definition}
Пусть $M$ "--- моноид, $a \in M$. Положим $a^0 = 1$ и $a^{e+1} = a^e \cdot a$ при $e \ge 0$.
Если $a$ обратим (например, $M$ "--- группа), то $a^{-e} = (a^{-1})^e$.
\end{definition}

\begin{proposition}\label{l6:powlaws}
Для $a \in M$ и $e, f \in \NN \cup \{0\}$: $a^{e+f} = a^e a^f$ и $(a^e)^f = a^{ef}$.
В частности, степени одного элемента перестановочны: $a^e a^f = a^f a^e$.
\end{proposition}

\begin{proof}
Первое равенство "--- индукция по $f$: при $f = 0$ это $a^e = a^e \cdot 1$, а
$a^{e+f+1} = a^{e+f} a = (a^e a^f) a = a^e (a^f a) = a^e a^{f+1}$ по ассоциативности.
Второе "--- индукция по $f$ с использованием первого: $(a^e)^{f+1} = (a^e)^f a^e = a^{ef} a^e
= a^{ef+e}$. Перестановочность: $a^e a^f = a^{e+f} = a^{f+e} = a^f a^e$.
\end{proof}

Именно эти два закона позволяют «собирать» показатель из частей: $a^{2e} = (a^e)^2$,
$a^{2e+1} = (a^e)^2 a$. Коммутативность самого моноида нигде не нужна.

\paragraph{Модель стоимости.} Мы считаем число \emph{умножений} в $M$, отдельно выделяя
\emph{возведения в квадрат} (у длинных чисел квадрат примерно вдвое дешевле произведения,
см. лекцию~2). Стоимость одного умножения по модулю $N$-значного $m$ "--- умножение
$N$-значных чисел плюс приведение по модулю (лекции~3 и~5), т.\,е. $\bigO(N^2)$ школьным
методом; поэтому возведение в степень $e < m$ стоит $\bigO(N^3)$ операций над цифрами.

Всюду ниже $e \ge 1$ и
\[
  e = (e_\lambda e_{\lambda-1} \dots e_1 e_0)_2 = \sum_{i=0}^{\lambda} e_i 2^i, \quad e_\lambda = 1,
  \qquad \lambda(e) = \floor{\log_2 e}, \qquad \nu(e) = \sum_i e_i,
\]
т.\,е. $\lambda(e) + 1$ "--- число двоичных разрядов, $\nu(e)$ "--- число единиц
(\emph{вес} показателя).

\section{Бинарный алгоритм}

\subsection{Справа налево}

Запишем $a^e = \prod_{i:\, e_i = 1} a^{2^i}$. Степени $a^{2^i}$ получаются последовательными
возведениями в квадрат, а нужные из них перемножаются.

\begin{algorithm}[бинарный, справа налево]\label{l6:rl}
\begin{algorithmic}[1]
\Require $a \in M$, $e \ge 1$.
\Ensure $a^e$.
\State $y \gets 1$, $x \gets a$, $f \gets e$
\While{$f > 0$}
  \If{$f$ нечётно} $y \gets y \cdot x$ \EndIf
  \State $f \gets \floor{f/2}$
  \If{$f > 0$} $x \gets x^2$ \EndIf
\EndWhile
\State \Return $y$
\end{algorithmic}
\end{algorithm}

\begin{theorem}\label{l6:rl-correct}
Алгоритм~\ref{l6:rl} возвращает $a^e$, выполняя $\lambda(e)$ возведений в квадрат и $\nu(e)$
умножений, из которых первое "--- на $y = 1$ "--- можно не выполнять.
\end{theorem}

\begin{proof}
\emph{Инвариант} цикла (в начале каждой итерации): $y \cdot x^f = a^e$. Вначале $y = 1$, $x = a$,
$f = e$. Если $f = 2f' + \delta$, $\delta \in \{0, 1\}$, то после итерации
$y' = y x^\delta$, $x' = x^2$, $f' = \floor{f/2}$, и
$y' (x')^{f'} = y x^{\delta} x^{2f'} = y x^f = a^e$ (законы
утверждения~\ref{l6:powlaws}). Значение $f$ строго убывает, а при $f = 0$ инвариант даёт
$y = a^e$. (Пропуск возведения в квадрат на последней итерации не нарушает рассуждение:
значение $x$ больше не используется.)

Подсчёт: цикл выполняется столько раз, сколько двоичных разрядов у $e$, т.\,е. $\lambda(e) + 1$
раз, и возведение в квадрат делается на всех итерациях, кроме последней. Умножение
выполняется на итерациях с нечётным $f$, т.\,е. при $e_i = 1$; таких $\nu(e)$.
\end{proof}

\begin{example}\label{l6:ex-rl}
$7^{45} \bmod 101$, $45 = 101101_2$. Столбцы: бит $e_i$, $x = 7^{2^i}$, $y$ после шага.
\[
\begin{array}{c|cccccc}
  i & 0 & 1 & 2 & 3 & 4 & 5 \\ \hline
  e_i & 1 & 0 & 1 & 1 & 0 & 1 \\
  x = 7^{2^i} & 7 & 49 & 78 & 24 & 71 & 92 \\
  y & 7 & 7 & 41 & 75 & 75 & 32
\end{array}
\]
Итак, $7^{45} \equiv 32 \pmod{101}$: $5$ возведений в квадрат и $3$ умножения
(не считая $y = 1 \cdot 7$).
\end{example}

\subsection{Слева направо}

Схема Горнера для $e = (\dots((e_\lambda \cdot 2 + e_{\lambda-1}) \cdot 2 + e_{\lambda-2})\dots)\cdot 2 + e_0$
даёт второй вариант: биты читаются от старшего к младшему.

\begin{algorithm}[бинарный, слева направо]\label{l6:lr}
\begin{algorithmic}[1]
\Require $a \in M$, $e = (e_\lambda \dots e_0)_2$, $e_\lambda = 1$.
\Ensure $a^e$.
\State $r \gets a$
\For{$i = \lambda - 1, \lambda - 2, \dots, 0$}
  \State $r \gets r^2$
  \If{$e_i = 1$} $r \gets r \cdot a$ \EndIf
\EndFor
\State \Return $r$
\end{algorithmic}
\end{algorithm}

\begin{theorem}\label{l6:lr-correct}
Алгоритм~\ref{l6:lr} возвращает $a^e$ и выполняет ровно $\lambda(e)$ возведений в квадрат
и $\nu(e) - 1$ умножений, всего $\lambda(e) + \nu(e) - 1 \le 2\lambda(e)$ операций.
\end{theorem}

\begin{proof}
Обозначим $E_i = (e_\lambda e_{\lambda-1} \dots e_i)_2 = \floor{e / 2^i}$ "--- число, образованное
старшими битами. \emph{Инвариант}: после шага с номером $i$ выполнено $r = a^{E_i}$. Перед циклом
$r = a = a^{E_\lambda}$, так как $E_\lambda = 1$. Шаг $i$ переводит $r = a^{E_{i+1}}$ в
$r^2 a^{e_i} = a^{2E_{i+1} + e_i} = a^{E_i}$. После шага $i = 0$ получаем $r = a^{E_0} = a^e$.

Возведение в квадрат выполняется на каждом из $\lambda$ шагов, умножение "--- на шагах с $e_i = 1$,
$i < \lambda$, т.\,е. $\nu(e) - 1$ раз. Так как $\nu(e) \le \lambda(e) + 1$, всего
не более $2\lambda(e)$ операций.
\end{proof}

\begin{corollary}\label{l6:average}
Для показателя, равномерно распределённого среди $(\lambda + 1)$-битовых чисел, среднее число
умножений бинарного алгоритма равно $\lambda/2$, а среднее общее число операций "--- $\frac32\lambda$.
\end{corollary}

\begin{proof}
Биты $e_0, \dots, e_{\lambda-1}$ независимы и равны $1$ с вероятностью $\frac12$, поэтому
$\mathbb{E}[\nu(e) - 1] = \lambda/2$.
\end{proof}

Тот же пример $7^{45} \bmod 101$ слева направо: после битов $1, 0, 1, 1, 0, 1$ значения $r$ равны
$7, 49, 41, 51, 76, 32$ (это $7^{1}, 7^{2}, 7^{5}, 7^{11}, 7^{22}, 7^{45}$); квадраты
$49, 78, 65, 76, 19$.

\taskexample{exp-binary}

\begin{remark}
Оба варианта выполняют одинаковое число операций, но различаются в деталях.
Слева направо умножают всегда на \emph{одно и то же} $a$: если $a$ мало (например, $a = 2$ в тесте
Ферма), умножение на $a$ "--- это сдвиг или умножение на короткое число, и оно почти бесплатно; кроме
того, только этот вариант обобщается на «окна» (разделы~\ref{l6:sec-qary}--\ref{l6:sec-sliding}).
Справа налево не нужно заранее знать длину показателя (биты можно получать «на лету»
делением на $2$), а возведения в квадрат не зависят от битов и могут выполняться параллельно
с умножениями.
\end{remark}

\checkyourself{Сколько умножений нужно для $a^{65537}$ (стандартный открытый показатель RSA,
$65537 = 2^{16} + 1$)? А для $a^{2^{16} - 1}$?}

\checkyourself{Покажите, что $\begin{pmatrix}1 & 1\\ 1 & 0\end{pmatrix}^n =
\begin{pmatrix}F_{n+1} & F_n\\ F_n & F_{n-1}\end{pmatrix}$, где $F_n$ "--- числа Фибоначчи.
Сколько умножений матриц $2 \times 2$ нужно, чтобы найти $F_{1000} \bmod m$?}

\section{Модификации бинарного алгоритма}

\subsection{$2^k$-арный метод}\label{l6:sec-qary}

Запишем показатель по основанию $q = 2^k$: $e = (d_{s} d_{s-1} \dots d_0)_q$, $0 \le d_j < q$,
$d_s \ne 0$. Двоичные цифры группируются по $k$: каждая $q$-ичная цифра "--- это $k$ битов.

\begin{algorithm}[$2^k$-арный, слева направо]\label{l6:qary}
\begin{algorithmic}[1]
\Require $a \in M$, $e = (d_s \dots d_0)_{2^k}$, $d_s \ne 0$.
\Ensure $a^e$.
\State $A_0 \gets 1$, $A_1 \gets a$; \textbf{для} $d = 2, \dots, 2^k - 1$: $A_d \gets A_{d-1} \cdot a$
  \Comment{предвычисление}
\State $r \gets A_{d_s}$
\For{$j = s - 1, \dots, 0$}
  \State $r \gets r^{2^k}$ \Comment{$k$ возведений в квадрат}
  \If{$d_j \ne 0$} $r \gets r \cdot A_{d_j}$ \EndIf
\EndFor
\State \Return $r$
\end{algorithmic}
\end{algorithm}

\begin{theorem}\label{l6:qary-count}
Алгоритм~\ref{l6:qary} возвращает $a^e$. Он выполняет $2^k - 2$ умножений на этапе
предвычисления, не более $\lambda(e)$ возведений в квадрат и не более
$s \le \lambda(e)/k$ умножений в основном цикле; в среднем умножений в основном
цикле $(1 - 2^{-k})\,s$.
\end{theorem}

\begin{proof}
Корректность "--- как в теореме~\ref{l6:lr-correct}: при $E_j = (d_s \dots d_j)_q$ инвариант
$r = a^{E_j}$ сохраняется, поскольку $E_j = q E_{j+1} + d_j$ и
$r^{q} A_{d_j} = a^{q E_{j+1} + d_j}$. Таблица $A_2, \dots, A_{q-1}$ требует $q - 2$ умножений.
Число цифр $s + 1 = \ceil{(\lambda + 1)/k}$, и так как $d_s \ge 1$, то
$q^s \le e$, т.\,е. $s \le \lambda/k$; возведений в квадрат $ks \le \lambda$, умножений не
больше $s$. Ненулевая цифра в случайном показателе встречается с вероятностью $1 - 2^{-k}$.
\end{proof}

Итого около $\lambda + \frac{\lambda}{k} + 2^k$ операций: с ростом $k$ число умножений в цикле
падает, а предвычисление растёт экспоненциально. Оптимальное $k$ при данном $\lambda$ легко найти
перебором (см. таблицу~\ref{l6:tab-k} ниже), а асимптотически получаем следующее.

\begin{theorem}[Брауэр, 1939]\label{l6:brauer}
$a^e$ можно вычислить за $\lambda(e) + (1 + o(1))\dfrac{\lambda(e)}{\log_2 \lambda(e)}$ умножений
при $e \to \infty$.
\end{theorem}

\begin{proof}
По теореме~\ref{l6:qary-count} общее число операций не больше
$\lambda + \lambda/k + 2^k$. Положим $L = \log_2 \lambda$ и $k = \floor{L - 2\log_2 L}$
(при больших $\lambda$ это $\ge 1$). Тогда $2^k \le 2^{L}/L^2 = \lambda / L^2 = o(\lambda / L)$,
а $k \ge L - 2\log_2 L - 1 = L\,(1 - o(1))$, откуда $\lambda/k \le (1 + o(1))\,\lambda/L$.
\end{proof}

Таким образом, «лишние» умножения сверх неизбежных $\lambda$ возведений в квадрат составляют
лишь $o(\lambda)$, тогда как у бинарного алгоритма их в среднем $\lambda/2$.

\begin{example}
$7^{45} \bmod 101$ при $k = 2$: $45 = (2\,3\,1)_4$, таблица $A_0, \dots, A_3 = 1, 7, 49, 40$.
Далее $r = A_2 = 49$; $r^4 = 24$, $r = 24 \cdot 40 = 51 = 7^{11}$; $r^4 = 19$,
$r = 19 \cdot 7 = 32 = 7^{45}$. Здесь $2$ умножения для таблицы, $4$ квадрата, $2$ умножения
в цикле, всего $8$ операций "--- столько же, сколько у бинарного метода: при коротком показателе
предвычисление не окупается.
\end{example}

\taskexample{exp-qary}

\subsection{Метод скользящего окна}\label{l6:sec-sliding}

Две идеи улучшают $2^k$-арный метод. Во-первых, нулевые биты между «окнами» можно
обрабатывать по одному (просто возводить в квадрат), а окна начинать только с единичного бита.
Во-вторых, окно можно всегда заканчивать единицей: если окно $w = (w' 0 \dots 0)_2$, то
$r^{2^s} a^w = \bigl(r^{2^{s - t}} a^{w'}\bigr)^{2^t}$, где $t$ "--- число нулей в конце.
Тогда все значения окон \emph{нечётны}, и предвычислять нужно вдвое меньше степеней:
$a, a^3, a^5, \dots, a^{2^k - 1}$.

\begin{algorithm}[скользящее окно ширины $k$]\label{l6:sliding}
\begin{algorithmic}[1]
\Require $a \in M$, $e = (e_\lambda \dots e_0)_2$, $k \ge 1$.
\Ensure $a^e$.
\State $A_1 \gets a$, $a_2 \gets a^2$; \textbf{для} $d = 3, 5, \dots, 2^k - 1$:
  $A_d \gets A_{d-2} \cdot a_2$ \Comment{предвычисление}
\State $r \gets 1$, $i \gets \lambda$
\While{$i \ge 0$}
  \If{$e_i = 0$}
    \State $r \gets r^2$, \quad $i \gets i - 1$
  \Else
    \State найти наименьшее $l \ge \max(i - k + 1, 0)$ с $e_l = 1$; \quad $w \gets (e_i \dots e_l)_2$
    \State $r \gets r^{2^{i - l + 1}} \cdot A_w$, \quad $i \gets l - 1$
      \Comment{первое окно: $r \gets A_w$}
  \EndIf
\EndWhile
\State \Return $r$
\end{algorithmic}
\end{algorithm}

\begin{theorem}\label{l6:sliding-count}
Алгоритм~\ref{l6:sliding} возвращает $a^e$. Пусть битовая запись $e$ разбита им на $t$ окон
с единицами и некоторое число одиночных нулей, а первое окно имеет длину $\ell_1$. Тогда
алгоритм выполняет $1 + (\lambda + 1 - \ell_1)$ возведений в квадрат и
$(2^{k-1} - 1) + (t - 1)$ умножений, причём $t \le \ceil{(\lambda + 1)/k}$.
\end{theorem}

\begin{proof}
Пусть перед очередным шагом обработаны биты $e_\lambda, \dots, e_{i+1}$, т.\,е.
\emph{инвариант} "--- $r = a^{E_{i+1}}$, $E_{i+1} = \floor{e/2^{i+1}}$ ($E_{\lambda+1} = 0$, $r = 1$).
Шаг с нулевым битом переводит его в $a^{2E_{i+1}} = a^{E_i}$. Шаг с окном $w = (e_i \dots e_l)_2$
длины $i - l + 1$ даёт $a^{2^{i-l+1} E_{i+1} + w} = a^{E_l}$, и инвариант сохраняется
при $i \gets l - 1$. Окно оканчивается единицей, поэтому $w$ нечётно и $A_w$ есть в таблице;
$w \le 2^k - 1$, так как длина окна не больше $k$. В конце $i = -1$ и $r = a^{E_0} = a^e$.

Подсчёт: каждый бит, кроме битов первого окна (там $r = 1$ и квадраты не нужны), вызывает ровно
одно возведение в квадрат; ещё один квадрат "--- $a_2$ в предвычислении. Умножения: $2^{k-1} - 1$
в таблице и по одному на каждое окно, кроме первого. Наконец, сопоставим окну, начатому
с бита $i$, отрезок битов $e_i, e_{i-1}, \dots, e_{\max(i-k+1, 0)}$. Окно занимает начало
отрезка, а остальные биты отрезка "--- нули, которые обрабатываются по одному; поэтому следующее
окно начинается за пределами отрезка. Отрезки не пересекаются и имеют длину $k$ (кроме,
быть может, последнего), откуда $t \le \ceil{(\lambda+1)/k}$.
\end{proof}

В среднем после окна идёт один нулевой бит (до следующей единицы), поэтому на одно окно
приходится около $k + 1$ битов, и число умножений в основном цикле близко к $\lambda/(k + 1)$
(точный анализ см. в~\cite[разд.~4.6.3]{knuth-taocp2}). Итак, при той же таблице
скользящее окно ширины $k$ ведёт себя как $2^{k+1}$-арный метод, но с вдвое меньшим
предвычислением.

\begin{table}[ht]
\centering
\caption{Оптимальная ширина окна и число операций в модели
$\lambda + 2^k - 2 + \frac{\lambda}{k}(1 - 2^{-k})$ ($2^k$-арный метод) и
$\lambda + 2^{k-1} + \frac{\lambda}{k+1}$ (скользящее окно); у бинарного метода в среднем
$1{,}5\lambda$.}\label{l6:tab-k}
\begin{tabular}{r|cc|cc|c}
  $\lambda$ & $k$ & $2^k$-арный & $k$ & окно & бинарный \\ \hline
  $160$ & $4$ & $212$ & $4$ & $200$ & $240$ \\
  $512$ & $5$ & $641$ & $5$ & $613$ & $768$ \\
  $1024$ & $5$ & $1252$ & $6$ & $1202$ & $1536$ \\
  $2048$ & $6$ & $2446$ & $7$ & $2368$ & $3072$
\end{tabular}
\end{table}

\begin{example}\label{l6:ex-sliding}
$e = 1903 = 11101101111_2$, $k = 3$: окна $111\,|\,0\,|\,11\,|\,0\,|\,111\,|\,1$, т.\,е.
$e = ((7 \cdot 2 \cdot 2^2 + 3) \cdot 2 \cdot 2^3 + 7) \cdot 2 + 1$. Для $a = 7$, $m = 101$:
таблица $a^2 = 49$, $a^3 = 40$, $a^5 = 41$, $a^7 = 90$; далее
\[
\begin{array}{c|l|c|r|r}
  \text{окно} & \text{квадраты} & \text{множитель} & r & \text{показатель} \\ \hline
  111 & \text{---} & a^7 & 90 & 7 \\
  0 & 20 & \text{---} & 20 & 14 \\
  11 & 97,\ 16 & a^3 & 34 & 59 \\
  0 & 45 & \text{---} & 45 & 118 \\
  111 & 5,\ 25,\ 19 & a^7 & 94 & 951 \\
  1 & 49 & a & 40 & 1903
\end{array}
\]
$7^{1903} \equiv 40 \pmod{101}$; $9$ квадратов и $6$ умножений против $10 + 8 = 18$ операций
бинарного метода. При $k = 2$ окна такие: $11\,|\,1\,|\,0\,|\,11\,|\,0\,|\,11\,|\,11$.
\end{example}

\taskexample{exp-sliding}

\checkyourself{Почему в методе скользящего окна предвычисляются только нечётные степени?
Сколько операций потребует метод окна ширины $3$ для $e = 2^{10}$ и сколько "--- бинарный
метод?}

\subsection{Знаковые цифры и NAF}\label{l6:sec-naf}

Если в $M$ обратный элемент вычисляется дёшево, полезны цифры $-1, 0, 1$: тогда
$r \gets r^2 a^{-1}$ стоит столько же, сколько $r \gets r^2 a$. Так обстоит дело в группе точек
эллиптической кривой, где $-P$ получается сменой знака одной координаты (лекция~9 второго семестра).
В $\Znx{m}$ обращение дорого (расширенный алгоритм Евклида, лекция~3), но $a^{-1}$ нужен один раз,
и знаковые цифры тоже могут окупиться.

Например, $2^n - 1 = (11\dots1)_2$ имеет вес $n$, а запись $2^n - 1 = 2^n - 2^0$ "--- всего две
ненулевые цифры.

\begin{definition}
Запись $e = \sum_{i} d_i 2^i$ с $d_i \in \{-1, 0, 1\}$ называется \emph{несмежной формой}
(non-adjacent form, NAF), если $d_i d_{i+1} = 0$ при всех $i$, т.\,е. никакие две ненулевые
цифры не стоят рядом.
\end{definition}

\begin{theorem}\label{l6:naf}
Каждое натуральное $e$ имеет единственную NAF. Её длина не больше $\lambda(e) + 2$.
\end{theorem}

\begin{proof}
\emph{Существование} "--- индукция по $e$. Если $e$ чётно, полагаем $d_0 = 0$ и берём NAF числа
$e/2$. Если $e$ нечётно, полагаем $d_0 = 2 - (e \bmod 4) \in \{1, -1\}$; тогда
$e - d_0 \equiv 0 \pmod 4$, т.\,е. $(e - d_0)/2$ чётно, и его NAF начинается с нуля: $d_1 = 0$.
Так как $(e - d_0)/2 < e$ при $e \ge 2$ (а $e = 1$ "--- это $d_0 = 1$), индукция проходит.

\emph{Единственность}. Пусть $\sum d_i 2^i = \sum d'_i 2^i$ "--- две NAF. Цифра $d_0$ определена
чётностью: $d_0 \equiv e \pmod 2$. Если $e$ нечётно, то $d_1 = 0$, поэтому
$e = d_0 + 4(\dots)$, т.\,е. $d_0 \equiv e \pmod 4$, и $d_0 \in \{1, -1\}$ определяется однозначно.
Дальше переходим к $(e - d_0)/2$.

\emph{Длина}. Пусть старшая ненулевая цифра NAF стоит в позиции $j$. Так как ненулевые цифры
не соседствуют, $\abs[\big]{\sum_{i<j} d_i 2^i} \le 2^{j-2} + 2^{j-4} + \cdots < 2^j/3$; поэтому
$d_j = 1$ (иначе $e < 0$) и $e > 2^j - 2^j/3 > 2^{j-1}$. Значит, $j - 1 < \log_2 e$, т.\,е.
$j \le \lambda(e) + 1$, и в NAF не больше $\lambda(e) + 2$ цифр.
\end{proof}

Известно, что NAF имеет наименьший вес среди всех записей с цифрами $\{-1, 0, 1\}$, а средняя доля
ненулевых цифр в ней равна $1/3$ (а не $1/2$, как у обычной двоичной записи)~\cite[гл.~11]{galbraith}.
Поэтому возведение в степень по NAF стоит в среднем $\lambda + \lambda/3$ операций.
Например, $1903 = 2^{11} - 2^7 - 2^4 - 1$: в NAF четыре ненулевые цифры против восьми единиц.

\begin{algorithm}[NAF]
\begin{algorithmic}[1]
\Require $e \ge 1$.
\Ensure цифры $d_0, d_1, \dots$ несмежной формы $e$.
\State $i \gets 0$
\While{$e > 0$}
  \If{$e$ нечётно} $d_i \gets 2 - (e \bmod 4)$, $e \gets e - d_i$ \Else\ $d_i \gets 0$ \EndIf
  \State $e \gets e/2$, \quad $i \gets i + 1$
\EndWhile
\end{algorithmic}
\end{algorithm}

\checkyourself{Найдите NAF чисел $45$ и $255$. Сколько операций нужно для $nP$ на эллиптической
кривой при $n = 255$ с NAF и без неё?}

\section{Аддитивные цепочки}

Все рассмотренные методы строят $a^e$ как последовательность степеней $a^{c_0}, a^{c_1}, \dots$,
каждая из которых "--- произведение двух предыдущих. Это приводит к комбинаторной задаче о
показателях.

\begin{definition}
\emph{Аддитивная цепочка} для $e$ "--- последовательность $1 = c_0 < c_1 < \cdots < c_L = e$,
в которой каждый член $c_i$, $i \ge 1$, равен сумме двух (возможно, совпадающих) предыдущих:
$c_i = c_j + c_k$, $0 \le k \le j < i$. Число $L$ "--- \emph{длина} цепочки; $\ell(e)$ "---
наименьшая длина цепочки для $e$.
\end{definition}

Каждая цепочка длины $L$ даёт способ вычислить $a^e$ за $L$ умножений: $a^{c_i} = a^{c_j} a^{c_k}$
(при $j = k$ "--- возведение в квадрат). Бинарный метод соответствует цепочке из удвоений и
прибавлений единицы, например для $45$:
$1, 2, 4, 5, 10, 11, 22, 44, 45$ (длина $\lambda + \nu - 1 = 8$).

\begin{theorem}\label{l6:chain-bounds}
$\ceil{\log_2 e} \le \ell(e) \le \lambda(e) + \nu(e) - 1$.
\end{theorem}

\begin{proof}
Верхняя оценка "--- цепочка бинарного метода (теорема~\ref{l6:lr-correct}). Нижняя: индукцией
по $i$ получаем $c_i \le 2^i$, так как $c_i = c_j + c_k \le 2 c_{i-1} \le 2^i$. Значит,
$e = c_L \le 2^L$, т.\,е. $L \ge \log_2 e$.
\end{proof}

Нижнюю оценку можно уточнить: $\ell(e) \ge \log_2 e + \log_2 \nu(e) - 2{,}13$ (Шёнхаге, 1975;
без доказательства, см.~\cite[разд.~4.6.3]{knuth-taocp2}). Верхняя оценка
теоремы~\ref{l6:brauer} асимптотически точна: для почти всех $e$ (П.~Эрдёш, 1960)
$\ell(e) = \lambda(e) + (1 + o(1))\,\lambda(e)/\log_2 \lambda(e)$.

\begin{example}
$\ell(15) = 5$: бинарный метод даёт цепочку $1, 2, 3, 6, 7, 14, 15$ длины $6$, а цепочки
$1, 2, 3, 6, 12, 15$ и $1, 2, 4, 5, 10, 15$ имеют длину $5$. Первая из них соответствует
\emph{методу множителей}: $15 = 3 \cdot 5$, сначала $b = a^3$ (две операции), затем
$b^5 = (b^2)^2 b$ (ещё три). Длина $4$ невозможна: из $15 = c_4 \le 2c_3$ следует $c_3 \ge 8$,
а так как $c_i \le 2^i$, то $c_3 = 8$ и цепочка "--- $1, 2, 4, 8$; но $15$ не является суммой
двух чисел из $\{1, 2, 4, 8\}$. Аналогично
$\ell(31) = 7 < 8$: $1, 2, 4, 8, 10, 20, 30, 31$.
\end{example}

\taskexample{exp-chain}

\begin{remark}
Задача нахождения кратчайшей аддитивной цепочки трудна. Для близкой задачи "--- кратчайшей
цепочки, содержащей \emph{все} числа заданного набора, "--- доказана NP-полнота
(Дауни, Леонг, Сети, 1981); для одного числа $e$ полиномиальный алгоритм не известен, а
перебор растёт экспоненциально (алгоритм, которым вычислены цепочки в задачах, "--- поиск
в глубину с ограничением длины и отсечением $c_i \cdot 2^{L - i} < e$). На практике кратчайшие
цепочки ищут один раз для \emph{фиксированных} показателей (например, $p - 2$ при обращении
по малой теореме Ферма в $\GF{p}$ для стандартного простого $p$), а для случайных показателей
используют окна. Открыта и гипотеза Шольца"--~Брауэра $\ell(2^n - 1) \le n - 1 + \ell(n)$.
\end{remark}

\checkyourself{Докажите, что $\ell(2e) \le \ell(e) + 1$ и $\ell(ef) \le \ell(e) + \ell(f)$.
Почему отсюда не следует $\ell(2e) = \ell(e) + 1$?}

\section{Возведение в степень методом Монтгомери}\label{l6:sec-montgomery}

\subsection{Напоминание: умножение Монтгомери}

В модулярном возведении в степень каждое из $\approx 1{,}5\lambda$ умножений сопровождается
приведением по модулю, т.\,е. делением $2N$-значного числа на $N$-значное $m$ (лекция~3).
Деление "--- самая медленная из арифметических операций: в школьном алгоритме очередную цифру
частного приходится \emph{угадывать и проверять} (trial division). Метод П.~Монтгомери
(1985,~\cite{montgomery}; статья так и называется "--- «Модулярное умножение без пробного
деления») заменяет деление на $m$ делением на степень основания системы счисления, т.\,е. сдвигом
(подробно "--- в лекции~5).

Пусть $m$ нечётно, $R = 2^k > m$ (на практике $R = B^n$, где $B = 2^{64}$ "--- основание),
$\gcd(R, m) = 1$, и заранее вычислены $m' = -m^{-1} \bmod R$ и $R^2 \bmod m$.

\begin{algorithm}[REDC]\label{l6:redc}
\begin{algorithmic}[1]
\Require $0 \le T < mR$.
\Ensure $\mathrm{REDC}(T) \equiv T R^{-1} \pmod m$, $0 \le \mathrm{REDC}(T) < m$.
\State $u \gets (T \bmod R) \cdot m' \bmod R$ \Comment{младшие $k$ битов, маска}
\State $t \gets (T + u m) / R$ \Comment{деление нацело, сдвиг на $k$ битов}
\If{$t \ge m$} $t \gets t - m$ \EndIf
\State \Return $t$
\end{algorithmic}
\end{algorithm}

\begin{lemma}\label{l6:redc-lemma}
При $0 \le T < mR$ алгоритм~\ref{l6:redc} возвращает единственное число $t \in [0, m)$ с
$t \equiv T R^{-1} \pmod m$.
\end{lemma}

\begin{proof}
$T + um \equiv T + T m' m \equiv T - T \equiv 0 \pmod R$, так как $m' m \equiv -1 \pmod R$;
поэтому деление на $R$ нацело. Далее $tR = T + um \equiv T \pmod m$, т.\,е. $t \equiv TR^{-1}$.
Наконец, $0 \le u < R$ и $T < mR$ дают $0 \le t < (mR + Rm)/R = 2m$, так что одного вычитания
достаточно.
\end{proof}

\begin{definition}
\emph{Форма Монтгомери} вычета $x \in \Zn{m}$ "--- число $\bar x = xR \bmod m$.
\end{definition}

Отображение $x \mapsto \bar x$ "--- биекция $\Zn{m} \to \Zn{m}$ (так как $R$ обратим по модулю
$m$), сохраняющая сложение, а умножение в формах Монтгомери выполняется через REDC:
\begin{equation}\label{l6:montmul}
  \mathrm{REDC}(\bar x\,\bar y) \equiv xR \cdot yR \cdot R^{-1} = (xy) R \equiv \overline{xy} \pmod m,
\end{equation}
причём $0 \le \bar x\bar y < m^2 < mR$, так что лемма~\ref{l6:redc-lemma} применима.
Переход в форму Монтгомери и обратно:
$\mathrm{REDC}\bigl(x \cdot (R^2 \bmod m)\bigr) = \bar x$ и $\mathrm{REDC}(\bar x) = x$.

\subsection{Алгоритм}

Если одно умножение по Монтгомери выполнять «как есть» "--- перевести сомножители в форму
Монтгомери, перемножить и перевести обратно, "--- то выигрыша нет: преобразования стоят
дороже одного деления. Но в возведении в степень \emph{все} промежуточные значения можно держать
в форме Монтгомери: перевод туда и обратно делается один раз на всё вычисление.

\begin{algorithm}[возведение в степень по Монтгомери]\label{l6:montexp}
\begin{algorithmic}[1]
\Require нечётный модуль $m$, $R = 2^k > m$, $m'$, $R^2 \bmod m$; $0 \le a < m$;
  $e = (e_\lambda \dots e_0)_2$.
\Ensure $a^e \bmod m$.
\State $\bar a \gets \mathrm{REDC}\bigl(a \cdot (R^2 \bmod m)\bigr)$ \Comment{$\bar a = aR \bmod m$}
\State $\bar x \gets \bar a$
\For{$i = \lambda - 1, \dots, 0$}
  \State $\bar x \gets \mathrm{REDC}(\bar x^2)$
  \If{$e_i = 1$} $\bar x \gets \mathrm{REDC}(\bar x \cdot \bar a)$ \EndIf
\EndFor
\State \Return $\mathrm{REDC}(\bar x)$ \Comment{выход из формы Монтгомери}
\end{algorithmic}
\end{algorithm}

\begin{theorem}\label{l6:montexp-correct}
Алгоритм~\ref{l6:montexp} возвращает $a^e \bmod m$. Он выполняет $\lambda(e) + \nu(e) - 1$
умножений с REDC и ещё два вызова REDC (вход и выход); ни одного деления на $m$ в основном
цикле нет.
\end{theorem}

\begin{proof}
\emph{Инвариант}: после шага $i$ выполнено $\bar x = \overline{a^{E_i}} = a^{E_i} R \bmod m$,
$0 \le \bar x < m$, где $E_i = \floor{e/2^i}$. Перед циклом $\bar x = \bar a = \overline{a^{1}}$
по лемме~\ref{l6:redc-lemma} ($0 \le a \cdot (R^2 \bmod m) < m^2 < mR$). На шаге $i$
по~\eqref{l6:montmul} $\mathrm{REDC}(\bar x^2) = \overline{a^{2E_{i+1}}}$ и, если $e_i = 1$,
$\mathrm{REDC}(\overline{a^{2E_{i+1}}} \cdot \bar a) = \overline{a^{2E_{i+1} + 1}}$, т.\,е.
в обоих случаях $\overline{a^{E_i}}$; все аргументы REDC меньше $m^2 < mR$, а результаты лежат
в $[0, m)$, так что лемма применима на каждом шаге. После шага $i = 0$ имеем
$\bar x = a^e R \bmod m$, и $\mathrm{REDC}(\bar x) \equiv a^e R R^{-1} = a^e$, $0 \le \cdot < m$.
Подсчёт операций "--- как в теореме~\ref{l6:lr-correct}.
\end{proof}

\paragraph{Стоимость.} REDC для $N$-словного $m$ ($R = B^N$) выполняется пословно за
$N^2 + N$ умножений слов (лекция~5), т.\,е. одно модулярное умножение стоит около $2N^2$ умножений
слов "--- примерно столько же, сколько школьное деление, но без угадывания цифр частного,
нормализации и ветвлений. Вся работа, требующая деления, "--- вычисление $m'$ (обращение по
модулю $2^{64}$, несколько умножений слов методом Ньютона) и $R^2 \bmod m$ "--- выполняется
\emph{один раз для модуля} и в RSA или Диффи"--~Хеллмане используется многократно. Вход в форму
Монтгомери и выход "--- два REDC на всё возведение в степень против $\approx 1{,}5\lambda$
умножений в цикле. Метод Монтгомери так же сочетается с окнами: таблица $\bar a, \overline{a^3},
\dots$ строится в форме Монтгомери.

\begin{example}\label{l6:ex-mont}
$7^{45} \bmod 101$: $R = 2^7 = 128$, $m' = -101^{-1} \bmod 128 = 19$, $R^2 \bmod 101 = 22$,
$45 = 101101_2$. Звёздочка "--- выполнено вычитание $t - m$.
\[
\begin{array}{l|r|r|r|r}
  & T & u & t & \mathrm{REDC}(T) \\ \hline
  \bar a = \mathrm{REDC}(7 \cdot 22) & 154 & 110 & 88 & 88 \\
  \text{бит } 0:\ \bar x^2 & 7744 & 64 & 111 & 10^{*} \\
  \text{бит } 1:\ \bar x^2 & 100 & 108 & 86 & 86 \\
  \phantom{\text{бит } 1:\ }\bar x \bar a & 7568 & 48 & 97 & 97 \\
  \text{бит } 1:\ \bar x^2 & 9409 & 83 & 139 & 38^{*} \\
  \phantom{\text{бит } 1:\ }\bar x \bar a & 3344 & 48 & 64 & 64 \\
  \text{бит } 0:\ \bar x^2 & 4096 & 0 & 32 & 32 \\
  \text{бит } 1:\ \bar x^2 & 1024 & 0 & 8 & 8 \\
  \phantom{\text{бит } 1:\ }\bar x \bar a & 704 & 64 & 56 & 56 \\
  \text{выход: } \mathrm{REDC}(\bar x) & 56 & 40 & 32 & 32
\end{array}
\]
Проверка инварианта: после первого шага $\bar x = 10 = 7^2 \cdot 128 \bmod 101$. Ответ
$7^{45} \equiv 32 \pmod{101}$ совпадает с примером~\ref{l6:ex-rl}.
\end{example}

\taskexample{exp-montgomery}

\begin{remark}[чётный модуль]
При чётном $m$ число $R = 2^k$ не обратимо по модулю $m$, и метод неприменим напрямую. Пишут
$m = 2^t s$ с нечётным $s$, вычисляют $a^e \bmod 2^t$ (приведение по модулю $2^t$ "--- маска)
и $a^e \bmod s$ методом Монтгомери, а затем объединяют результаты по китайской теореме об остатках
(лекция~5; подробно "--- лекция~2 второго семестра). В криптографии
модули нечётны (простые $p$ или $N = pq$), так что этот случай редок, но на семинаре он встретится.
\end{remark}

\checkyourself{Почему в теореме~\ref{l6:montexp-correct} аргументы всех вызовов REDC меньше
$mR$? Что сломается, если выбрать $R < m$?}

\checkyourself{Как изменится алгоритм~\ref{l6:montexp}, если $a \equiv 0 \pmod m$ или
$\gcd(a, m) > 1$? Нужно ли, чтобы $a$ было обратимо по модулю $m$?}

\section{Атаки по побочным каналам и лестница Монтгомери}

В бинарном алгоритме умножение выполняется только при $e_i = 1$. Если показатель секретный
(расшифрование RSA, ключ Диффи"--~Хеллмана), то время работы и энергопотребление зависят от
секрета: П.~Кохер~\cite{kocher-timing} (1996) показал, как по времени выполнения многих
операций восстановить секретный показатель бит за битом, а по осциллограмме
энергопотребления смарт-карты последовательность «квадрат, квадрат, умножение, \dots»
читается непосредственно (\emph{простой анализ энергопотребления}, SPA). Признаки такого
рода "--- \emph{побочные каналы}. Защита "--- алгоритмы, у которых последовательность операций
не зависит от битов показателя (\emph{постоянное время}, constant time).

\begin{algorithm}[лестница Монтгомери]\label{l6:ladder}
\begin{algorithmic}[1]
\Require $a \in M$, $e = (e_\lambda \dots e_0)_2$.
\Ensure $a^e$.
\State $R_0 \gets 1$, $R_1 \gets a$
\For{$i = \lambda, \lambda - 1, \dots, 0$}
  \If{$e_i = 1$} $R_0 \gets R_0 R_1$, \quad $R_1 \gets R_1^2$
  \Else\ $R_1 \gets R_0 R_1$, \quad $R_0 \gets R_0^2$ \EndIf
\EndFor
\State \Return $R_0$
\end{algorithmic}
\end{algorithm}

\begin{theorem}\label{l6:ladder-correct}
Лестница Монтгомери возвращает $a^e$; на каждом бите она выполняет ровно одно умножение и одно
возведение в квадрат.
\end{theorem}

\begin{proof}
\emph{Инвариант}: после обработки бита $i$ выполнено $R_0 = a^{E_i}$, $R_1 = a^{E_i + 1}$
(до цикла $R_0 = a^0$, $R_1 = a^1$, что соответствует $E_{\lambda+1} = 0$). Если $e_i = 1$, то
$E_i = 2E_{i+1} + 1$, и $R_0 R_1 = a^{2E_{i+1} + 1} = a^{E_i}$, $R_1^2 = a^{2E_{i+1} + 2} = a^{E_i+1}$.
Если $e_i = 0$, то $E_i = 2E_{i+1}$, $R_0^2 = a^{E_i}$, $R_0 R_1 = a^{E_i + 1}$. В конце
$R_0 = a^{E_0} = a^e$. (Умножаются только степени $a$, поэтому коммутативность $M$ не нужна:
утверждение~\ref{l6:powlaws}.)
\end{proof}

Для $7^{45} \bmod 101$ пары $(R_0, R_1)$ после битов $1, 0, 1, 1, 0, 1$:
\[
  (7, 49),\ (49, 40),\ (41, 85),\ (51, 54),\ (76, 27),\ (32, 22).
\]

Ветвление по биту в реальной реализации заменяют перестановкой $(R_0, R_1)$ без условных переходов
(«constant-time swap» по маске). Лестница стоит $2(\lambda + 1)$ операций "--- дороже бинарного
метода в среднем на треть, "--- зато хорошо параллелится (умножение и квадрат на каждом шаге
независимы) и в группе точек эллиптической кривой позволяет вести вычисления только с
$x$-координатами (так устроен протокол X25519). Другие меры защиты: окна фиксированной
ширины (умножение выполняется и на нулевую цифру, $A_0 = 1$), \emph{ослепление}
показателя $e + r\cdot\ord(a)$ или основания со случайным $r$. Кроме того, сами операции должны
выполняться за постоянное время: например, финальное вычитание в REDC (строка~3
алгоритма~\ref{l6:redc}) тоже утечка, и его делают безусловным.

\checkyourself{Сколько операций выполнит лестница Монтгомери для $e = 2^{10}$ и для
$e = 2^{11} - 1$? Почему это хорошо с точки зрения защиты?}

\section{Итоги}

\begin{itemize}
  \item Бинарный алгоритм: $\lambda(e)$ квадратов и $\nu(e) - 1$ умножений, в среднем
    $1{,}5\lambda$ операций, в худшем случае $2\lambda$.
  \item Окна ($2^k$-арный метод, скользящее окно): $\lambda + \bigO(\lambda/k) + \bigO(2^k)$;
    при оптимальном $k$ "--- $\lambda + (1 + o(1))\lambda/\log_2\lambda$ (Брауэр), что
    асимптотически оптимально для почти всех $e$.
  \item Нижняя оценка любого метода, основанного только на умножениях, "--- $\ceil{\log_2 e}$.
  \item Знаковые цифры (NAF) уменьшают долю умножений до $1/3$, если обращение дёшево (эллиптические
    кривые).
  \item Модулярное возведение в степень выполняют в форме Монтгомери: одно преобразование
    на входе и одно на выходе, ни одного деления на $m$ внутри цикла.
  \item При секретном показателе используют алгоритмы постоянного времени (лестница Монтгомери).
\end{itemize}

\checkyourselfsection

\begin{furtherreading}
  \readitem[разд.~4.6.3]{knuth-taocp2}{бинарный и $2^k$-арный методы, аддитивные цепочки,
    теоремы Брауэра, Шёнхаге и Эрдёша, гипотеза Шольца"--~Брауэра}
  \readitem[гл.~6]{welschenbach}{модулярное возведение в степень: бинарный и $M$-арный методы,
    аддитивные цепочки и окна, возведение в степень по Монтгомери; реализация на Си}
  \readitem[разд.~9.2.1, 9.3]{crandall-pomerance}{метод Монтгомери; «лестницы» возведения
    в степень и их улучшения}
  \readitem[разд.~14.6]{hac}{алгоритмы 14.76--14.94: бинарные, $k$-арный, скользящее окно,
    возведение в степень по Монтгомери; аддитивные цепочки, фиксированный показатель и основание}
  \readitem[гл.~11]{galbraith}{возведение в степень в группах, знаковые представления и NAF}
  \readitem[разд.~7.4, 7.5, 7.8]{paar-pelzl}{square-and-multiply, ускорение RSA,
    атаки по побочным каналам}
  \readitem{montgomery}{оригинальная статья П.~Монтгомери}
  \readitem{kocher-timing}{атаки по времени выполнения}
  \readitem{joye-yen-ladder}{лестница Монтгомери и её свойства}
  \readitem{cp-algorithms}{раздел «Binary Exponentiation» с примерами применения}
  \readitem{paar-lectures}{видеолекция 12 о реализации RSA и square-and-multiply}
\end{furtherreading}
\end{lecture}
